\clearpage\pagestyle{otMain}\setlength{\otBlocEspacement}{2.5mm}\fontsize{11}{15}\selectfont\corrigesChapitre{21}{Instructions, répétitions et chemins}
\begin{blocLivre}{cor-21-1}\begin{corrige}{1}{Instruction ou séquence ?}\demarche{Compter les ordres élémentaires.}
Instructions : $\res{\text{« avancer de 3 cases » et « afficher le résultat »}}$. Séquences : les deux autres, formées de plusieurs instructions dans un ordre donné.
\end{corrige}\end{blocLivre}
\begin{blocLivre}{cor-21-2}\begin{corrige}{2}{Exécuter un programme}\demarche{Une ligne par instruction.}
$7\times2=14$ ; $14-3=11$ ; $11\times10=110$. Sortie : $\res{110}$.
\label{fig-cor-21-2}\figureGeo[corrige]{ins-2}
\end{corrige}\end{blocLivre}
\begin{blocLivre}{cor-21-3}\begin{corrige}{3}{Entrée et sortie}\demarche{Repérer ce qui est donné et ce qui est affiché.}
Entrée : $\res{\text{le nombre d’oranges}}$. Sortie : $\res{\text{la quantité de jus}}$. Pour 6 oranges : $6\times8=\res{48\ \mathrm{cL}}$.
\end{corrige}\end{blocLivre}
\begin{blocLivre}{cor-21-4}\begin{corrige}{4}{Suivre un chemin}\demarche{Suivre chaque instruction en notant l’orientation.}
Le trajet est tracé en vert : le robot arrive 5 cases à droite et 2 cases au-dessus de D, sur la case $\res{\text{R}}$.
\label{fig-cor-21-4}\figureGeo[corrige]{ins-4}
\end{corrige}\end{blocLivre}
\begin{blocLivre}{cor-21-5}\begin{corrige}{5}{L’ordre compte}\demarche{Exécuter dans l’ordre écrit.}
Premier : $(5+3)\times2=\res{16}$. Second : $5\times2+3=\res{13}$. Les mêmes instructions dans un autre ordre donnent une autre sortie.
\end{corrige}\end{blocLivre}
\begin{blocLivre}{cor-21-6}\begin{corrige}{6}{Répéter à la main}\demarche{Noter position et orientation à chaque tour.}
Tour 1 : 2 cases à droite, puis il se tourne vers le bas. Tour 2 : 2 cases vers le bas, puis vers la gauche. Tour 3 : 2 cases à gauche, puis vers le haut. Il arrive $\res{2\text{ cases sous le départ}}$, tourné $\res{\text{vers le haut}}$ (trajet en vert).
\label{fig-cor-21-6}\figureGeo[corrige]{ins-6}
\end{corrige}\end{blocLivre}
\begin{blocLivre}{cor-21-7}\begin{corrige}{7}{Écrire avec une boucle}\demarche{Repérer le morceau qui revient.}
« $\res{\text{Répéter 4 fois : avancer de 2 cases, tourner à gauche}}$ ». Le personnage trace un carré de 2 cases de côté et revient à son point de départ.
\end{corrige}\end{blocLivre}
\begin{blocLivre}{cor-21-8}\begin{corrige}{8}{Coder un chemin}\demarche{Découper le trajet en lignes droites.}
Par exemple : avancer de 2 cases ; tourner à gauche ; avancer de 3 cases ; tourner à droite ; avancer de 3 cases. Ce chemin est tracé en vert ; d’autres chemins conviennent.
\label{fig-cor-21-8}\figureGeo[corrige]{ins-8}
\end{corrige}\end{blocLivre}
\begin{blocLivre}{cor-21-9}\begin{corrige}{9}{Remonter un programme}\demarche{Défaire les instructions à l’envers.}
On remonte : $50-2=48$, puis $48:4=\res{12}$. Contrôle : $12\times4+2=50$.
\end{corrige}\end{blocLivre}
\begin{blocLivre}{cor-21-10}\begin{corrige}{10}{Le motif répété}\demarche{Chaque tour de boucle ajoute une marche.}
Il trace un $\res{\text{escalier}}$ de 4 marches, en vert. Il arrive 4 cases à droite et 4 cases au-dessus du départ, tourné vers la droite.
\label{fig-cor-21-10}\figureGeo[corrige]{ins-10}
\end{corrige}\end{blocLivre}
\begin{blocLivre}{cor-21-11}\begin{corrige}{11}{Le robot livreur}\demarche{Exécuter pas à pas.}
Positions successives : 2 cases à droite, puis 3 cases vers le haut, puis 4 cases à droite, puis 1 case vers le bas. Il s’arrête devant la maison $\res{\text{B}}$ (trajet en vert).
\label{fig-cor-21-11}\figureGeo[corrige]{ins-11}
\end{corrige}\end{blocLivre}
\begin{blocLivre}{cor-21-12}\begin{corrige}{12}{Le polygone}\demarche{Additionner les angles de rotation.}
Il tourne en tout de $6\times60=360^\circ$ et revient à son point de départ : il trace un $\res{\text{hexagone}}$ de 6 côtés égaux, en vert. Périmètre : $6\times3=\res{18\ \mathrm{cm}}$.
\label{fig-cor-21-12}\figureGeo[corrige]{ins-12}
\end{corrige}\end{blocLivre}
\begin{blocLivre}{cor-21-13}\begin{corrige}{13}{Le tableur}\demarche{Compter les pas de 4 depuis A1.}
Chaque cellule vaut 4 de plus que la précédente : A10 vaut $5+9\times4=\res{41}$. Pour 101 : $101-5=96$ et $96:4=24$ pas après A1, donc en $\res{\text{A25}}$.
\end{corrige}\end{blocLivre}
\begin{blocLivre}{cor-21-14}\begin{corrige}{14}{Trouver l’erreur}\demarche{Comparer le tracé à la figure voulue.}
Le tracé n’a que trois côtés : le carré n’est pas fermé. Il faut $\res{\text{répéter 4 fois}}$ : avancer de 4 cases, tourner à gauche. Le côté manquant est en vert.
\label{fig-cor-21-14}\figureGeo[corrige]{ins-14}
\end{corrige}\end{blocLivre}
\begin{blocLivre}{cor-21-15}\begin{corrige}{15}{Le plus court chemin}\demarche{Compter les déplacements indispensables.}
A est 5 cases à droite et 5 cases plus haut : il faut au moins $5+5=\res{10}$ cases. Le programme « répéter 5 fois : avancer de 1, tourner à gauche, avancer de 1, tourner à droite » réalise ce minimum en évitant les obstacles (en vert).
\label{fig-cor-21-15}\figureGeo[corrige]{ins-15}
\end{corrige}\end{blocLivre}
